\documentclass[10pt,a4paper]{amsart}
\usepackage{amsmath,amssymb,amsthm,graphicx}
\usepackage{verbatim}
\newenvironment{reference}{\comment}{\endcomment}
%\newenvironment{reference}{}{}
\newenvironment{slogan}{\comment}{\endcomment}
\newenvironment{history}{\comment}{\endcomment}

% For commutative diagrams we use Xy-pic
\usepackage[all]{xy}

% We use 2cell for 2-commutative diagrams.
\xyoption{2cell}
% Theorem environments.
%
\theoremstyle{plain}
\newtheorem{theorem}[subsection]{Theorem}
\newtheorem{proposition}[subsection]{Proposition}
\newtheorem{lemma}[subsection]{Lemma}

\theoremstyle{definition}
\newtheorem{definition}[subsection]{Definition}
\newtheorem{example}[subsection]{Example}
\newtheorem{exercise}[subsection]{Exercise}
\newtheorem{situation}[subsection]{Situation}

\theoremstyle{remark}
\newtheorem{remark}[subsection]{Remark}
\newtheorem{remarks}[subsection]{Remarks}

\numberwithin{equation}{subsection}

% Macros
%
\def\lim{\mathop{\mathrm{lim}}\nolimits}
\def\colim{\mathop{\mathrm{colim}}\nolimits}
\def\Spec{\mathop{\mathrm{Spec}}}
\def\Hom{\mathop{\mathrm{Hom}}\nolimits}
\def\Ext{\mathop{\mathrm{Ext}}\nolimits}
\def\SheafHom{\mathop{\mathcal{H}\!\mathit{om}}\nolimits}
\def\SheafExt{\mathop{\mathcal{E}\!\mathit{xt}}\nolimits}
\def\Sch{\mathit{Sch}}
\def\Mor{\mathop{\mathrm{Mor}}\nolimits}
\def\Ob{\mathop{\mathrm{Ob}}\nolimits}
\def\Sh{\mathop{\mathit{Sh}}\nolimits}
\def\NL{\mathop{N\!L}\nolimits}
\def\CH{\mathop{\mathrm{CH}}\nolimits}
\def\proetale{{pro\text{-}\acute{e}tale}}
\def\etale{{\acute{e}tale}}
\def\QCoh{\mathit{QCoh}}
\def\Ker{\mathop{\mathrm{Ker}}}
\def\Im{\mathop{\mathrm{Im}}}
\def\Coker{\mathop{\mathrm{Coker}}}
\def\Coim{\mathop{\mathrm{Coim}}}

% Boxtimes
%
\DeclareMathSymbol{\boxtimes}{\mathbin}{AMSa}{"02}

%
% Macros for moduli stacks/spaces
%
\def\QCohstack{\mathcal{QC}\!\mathit{oh}}
\def\Cohstack{\mathcal{C}\!\mathit{oh}}
\def\Spacesstack{\mathcal{S}\!\mathit{paces}}
\def\Quotfunctor{\mathrm{Quot}}
\def\Hilbfunctor{\mathrm{Hilb}}
\def\Curvesstack{\mathcal{C}\!\mathit{urves}}
\def\Polarizedstack{\mathcal{P}\!\mathit{olarized}}
\def\Complexesstack{\mathcal{C}\!\mathit{omplexes}}
% \Pic is the operator that assigns to X its picard group, usage \Pic(X)
% \Picardstack_{X/B} denotes the Picard stack of X over B
% \Picardfunctor_{X/B} denotes the Picard functor of X over B
\def\Pic{\mathop{\mathrm{Pic}}\nolimits}
\def\Picardstack{\mathcal{P}\!\mathit{ic}}
\def\Picardfunctor{\mathrm{Pic}}
\def\Deformationcategory{\mathcal{D}\!\mathit{ef}}
\usepackage{cleveref}
\crefname{originalreference}{original reference}{original references}
\Crefname{originalreference}{Original reference}{Original references}
\makeatletter
\newlabel{situation-setup}{{02QL}{source}{}{sourcepointer-1}{}}
\newlabel{situation-setup@cref}{{[originalreference][0][]02QL}{[1][0][]source}{}{}{}}
\makeatother
\crefname{originallemma}{original lemma}{original lemmas}
\Crefname{originallemma}{Original lemma}{Original lemmas}
\makeatletter
\newlabel{lemma-flat-pullback-bivariant}{{0B78}{source}{}{sourcepointer-2}{}}
\newlabel{lemma-flat-pullback-bivariant@cref}{{[originallemma][0][]0B78}{[1][0][]source}{}{}{}}
\makeatother
\crefname{originallemma}{original lemma}{original lemmas}
\Crefname{originallemma}{Original lemma}{Original lemmas}
\makeatletter
\newlabel{divisors-lemma-section-smooth-regular-immersion}{{067R}{source}{}{sourcepointer-3}{}}
\newlabel{divisors-lemma-section-smooth-regular-immersion@cref}{{[originallemma][0][]067R}{[1][0][]source}{}{}{}}
\makeatother
\crefname{originallemma}{original lemma}{original lemmas}
\Crefname{originallemma}{Original lemma}{Original lemmas}
\makeatletter
\newlabel{lemma-gysin-fundamental}{{0FBM}{source}{}{sourcepointer-4}{}}
\newlabel{lemma-gysin-fundamental@cref}{{[originallemma][0][]0FBM}{[1][0][]source}{}{}{}}
\makeatother
\title{A section of a smooth morphism is a regular immersion and is a right inverse of the flat pullback bivariant class}
\author{The Stacks Project Authors}
\date{stacks-project:chow.tex:local-snapshot-20261012; GNU FDL 1.2 or later}
\begin{document}
\maketitle
\noindent\textit{Editorial note.} All delivered mathematical text is the author's own, copied from the pinned Stacks Project chapter file chow.tex; no statement or proof step was rewritten, reordered or removed.\par\medskip
\noindent\textit{Editorial note.} Preamble provenance: the pinned chapter file opens with an input of the shared preamble and carries no preamble text of its own, so the project's single Stacks-wide preamble staged at sources/stacks-chow/preamble.tex is excerpted verbatim as two spans, lines 14-18 (Xy-pic setup) and lines 154-217 (the native theorem declarations and mathematical macros). Nothing was added, removed or redefined; the wrapper class and the amsmath/amssymb/amsthm/graphicx base packages are supplied by the assembler.\par\medskip
\noindent\textit{Editorial note.} The counted statement and proof contain no cross reference, no citation and no included file, so the delivered body is self contained.\par\medskip
\noindent\textit{Editorial note.} Canonical location of the counted proposition: Stacks tag 0FF1 (sources/stacks-chow/pinned-original-source.tex line 13435).\par\medskip
\noindent\textit{Allowed standard prerequisites (editorial).}\begin{itemize}
\item A section of a smooth morphism is a regular immersion. Divisors, Lemma labelled divisors-lemma-section-smooth-regular-immersion (Stacks tag 067R), invoked at proof lines 13445-13446 and again at 13454-13456. Supplies both the first assertion of the statement and the regular embeddedness of the intersection used in the cycle computation.
\item Flat pullback as a bivariant class. Same chapter, Lemma labelled lemma-flat-pullback-bivariant (source line 6131), invoked in the statement. Supplies the class p\textasciicircum{}* whose composite with s\textasciicircum{}! is asserted to be the identity.
\item The fundamental-class identity for a regularly embedded closed subscheme. Same chapter, Lemma labelled lemma-gysin-fundamental (source line 12076), invoked at proof line 13457. Supplies the conclusion s\textasciicircum{}! cap p\textasciicircum{}*[Z] = [Z] once the intersection is identified as a codimension r regular embedding.
\item The base situation (S, delta) for operational Chow groups. Same chapter, Situation labelled situation-setup (source line 1372), invoked in the statement. Fixes the dimension function and the augmented ring A\textasciicircum{}*(X) in which the identity is stated.
\end{itemize}\medskip
\ifdefined\hypertarget\hypertarget{sourcepointer-1}{}\fi
\noindent\textit{Source reference.} Original reference 02QL: sources/stacks-chow/pinned-original-source.tex, source lines 1372--1372. Same-chapter pointer into the author environment situation at source line 1372; invoked at proof line None of the counted proof. It is retained only as a navigation pointer; every load-bearing step is inside the delivered span.\par\medskip
\ifdefined\hypertarget\hypertarget{sourcepointer-2}{}\fi
\noindent\textit{Source reference.} Original lemma 0B78: sources/stacks-chow/pinned-original-source.tex, source lines 6131--6131. Same-chapter pointer into the author environment lemma at source line 6131; invoked at proof line None of the counted proof. It is retained only as a navigation pointer; every load-bearing step is inside the delivered span.\par\medskip
\ifdefined\hypertarget\hypertarget{sourcepointer-3}{}\fi
\noindent\textit{Source reference.} Original lemma 067R: sources/stacks-chow/tags.txt, source lines 7980--7980. Cross-chapter standard prerequisite located in the project tag list staged in the local snapshot; invoked at proof line 13446 of the counted proof. It is retained only as a navigation pointer; every load-bearing step is inside the delivered span.\par\medskip
\ifdefined\hypertarget\hypertarget{sourcepointer-4}{}\fi
\noindent\textit{Source reference.} Original lemma 0FBM: sources/stacks-chow/pinned-original-source.tex, source lines 12076--12076. Same-chapter pointer into the author environment lemma at source line 12076; invoked at proof line 13457 of the counted proof. It is retained only as a navigation pointer; every load-bearing step is inside the delivered span.\par\medskip
\begin{lemma}
\label{lemma-section-smooth}
Let $(S, \delta)$ be as in Situation \ref{situation-setup}.
Let $p : P \to X$ be a smooth morphism of schemes locally of finite type
over $S$ and let $s : X \to P$ be a section. Then $s$ is a
regular immersion and $1 = s^! \circ p^*$ in $A^*(X)^\wedge$
where $p^* \in A^*(P \to X)^\wedge$ is the bivariant class
of Lemma \ref{lemma-flat-pullback-bivariant}.
\end{lemma}

\begin{proof}
The first statement is Divisors, Lemma
\ref{divisors-lemma-section-smooth-regular-immersion}.
It suffices to show that $s^! \cap p^*[Z] = [Z]$ in
$\CH_*(X)$ for any integral closed subscheme $Z \subset X$
as the assumptions are preserved by base change by $X' \to X$
locally of finite type. After replacing $P$ by an open neighbourhood
of $s(Z)$ we may assume $P \to X$ is smooth of fixed relative dimension $r$.
Say $\dim_\delta(Z) = n$. Then every irreducible component of $p^{-1}(Z)$
has dimension $r + n$ and $p^*[Z]$ is given by $[p^{-1}(Z)]_{n + r}$.
Observe that $s(X) \cap p^{-1}(Z) = s(Z)$ scheme theoretically. Hence by the
same reference as used above $s(X) \cap p^{-1}(Z)$ is a closed subscheme
regularly embedded in $\overline{p}^{-1}(Z)$ of codimension $r$.
We conclude by Lemma \ref{lemma-gysin-fundamental}.
\end{proof}

\end{document}
